\section{Matrix geometry in arbitrary input dimension}
\label{sec:matrixmetric76}

The binary effect body has a finite-use description that does not
require eigenvector coordinates.  Its regularization is the same at
every spectral rank, and two-dimensional compressions supply lower
witnesses for the full matrix problem.  Throughout this section
$d,N\ge1$ are integers and
\begin{equation}\label{eq:matrixbody76}
 \mathfrak E_d=\{E\in\End(\C^d):E=E^*,\ 0\preceq E\preceq I_d\}.
\end{equation}
The measurement $\mathcal M_E$ has the ordered effects $E,I_d-E$
and no residual quantum output.  The distance $d_N$ has the same
reference-assisted adaptive meaning as in
Theorem~\ref{thm:biasedmetric75}.  Write $d_N^{\rm na}$ for its
restriction to nonadaptive experiments, allowing entanglement across
the calls and a retained reference.  These experiments may depend on
the known pair of hypotheses.

Put $\tau=N^{-1}$ and $V_E=E(I_d-E)$.  On the real Hilbert space
of Hermitian matrices with Hilbert--Schmidt inner product, define
\begin{align}
 \mathcal A_E(H)&=V_EH+HV_E+\tau H,
                                      \label{eq:matrixoperator76}\\
 g_{N,E}(H)^2&=N\ip{H}{\mathcal A_E^{-1}(H)}_{\rm HS}.
                                      \label{eq:matrixform76}
\end{align}
The operator $\mathcal A_E$ is strictly positive, including at a
singular effect.  For a pair $E,F$ set
\begin{equation}\label{eq:matrixmodulus76}
 M=\tfrac12(E+F),\qquad H=F-E,\qquad
 Q_N(E,F)=g_{N,M}(H).
\end{equation}
In any orthonormal eigenbasis of $M$, with eigenvalues $m_i$ and
$v_i=m_i(1-m_i)$, this is
\begin{equation}\label{eq:matrixcoordinates76}
 Q_N(E,F)^2=N\sum_{i,j=1}^d
                \frac{|H_{ij}|^2}{v_i+v_j+N^{-1}}.
\end{equation}
Definition~\eqref{eq:matrixmodulus76} makes the expression independent
of all choices at repeated eigenvalues.  We use it as a comparison
modulus; a triangle inequality for $Q_N$ is not required.

The inverse Sylvester operator is familiar from the Bures metric
$g^{\rm B}$ on the positive definite cone \cite{Dittmann76}.
Its normalization gives the algebraic identity
\[
 g_{N,E}(H)^2=2N\,g^{\rm B}_{V_E+\tau I_d/2}(H,H).
\]
Here $H$ remains the effect tangent; the differential of the
variance map $E\mapsto V_E$ is $H-EH-HE$.

\begin{theorem}[Finite-use geometry of the matrix effect body]
\label{thm:matrixmetric76}
For every $E,F\in\mathfrak E_d$,
\begin{equation}\label{eq:matrixmetric76}
 \frac1{8192d}\min\{1,Q_N(E,F)\}
 \le d_N^{\rm na}(\mathcal M_E,\mathcal M_F)
 \le d_N(\mathcal M_E,\mathcal M_F)
 \le\min\{2,8Q_N(E,F)\}.
\end{equation}
The constants are uniform in the horizon and in both effects,
without a spectral margin or a condition on eigenvalue multiplicities.
The adaptive upper bound permits arbitrary retained references,
common feedback, and bounded public stopping.  Each lower witness
uses inputs in one fixed subspace of dimension at most two.  Moreover,
\begin{equation}\label{eq:matrixoneuse76}
 d_1(\mathcal M_E,\mathcal M_F)=2\|E-F\|_{\rm op}.
\end{equation}
\end{theorem}

The theorem includes every rank of projection and all of the
intermediate singular strata.  The variance $E(I_d-E)$, rather than
the smallest eigenvalue of either outcome, determines the matrix
weights.  The factor $N^{-1}$ makes the formula finite on the entire
closed body.  The proof first establishes a path inequality valid for
every common adaptive tester, and then obtains the converse from
one- and two-dimensional experiments.

\subsection{A horizontal measurement dilation}

Kraus-gauge bounds for channel estimation were developed in
\cite{FI76,DKG67,DM76}.  Their integration along channel paths
appears in \cite[Section~III, equation~(26)]{YF76} and, with
general adaptive bounds, in \cite[Section~IV.A]{SD76}.  We use an
explicit horizontal lift for binary effects and control the
finite-horizon regularization by a trace-norm remainder.

\begin{lemma}[Horizontal derivative and regularized path bound]
\label{lem:matrixpath76}
Every continuously differentiable path $G:[0,1]\to\mathfrak E_d$
satisfies
\begin{equation}\label{eq:matrixpath76}
 d_N(\mathcal M_{G(0)},\mathcal M_{G(1)})
       \le4\int_0^1 g_{N,G(t)}(G'(t))\,dt.
\end{equation}
\end{lemma}
\begin{proof}
We begin with an interior effect $G$, so that $G$ and $I_d-G$
are invertible.  Set $S=\sqrt{G(I_d-G)}$.  For a Hermitian matrix
$K$ put $H_0=SK+KS$.  At the two factors
$A_1=\sqrt G$, $A_0=\sqrt{I_d-G}$ prescribe derivatives
\begin{equation}\label{eq:horizontalderivatives76}
 B_1=\sqrt{I_d-G}\,K,\qquad B_0=-\sqrt G\,K.
\end{equation}
They obey
\begin{align}
 A_1^*B_1+B_1^*A_1&=H_0,&
 A_0^*B_0+B_0^*A_0&=-H_0,\notag\\
 \sum_{y=0}^1 A_y^*B_y&=0,&
 \sum_{y=0}^1 B_y^*B_y&=K^2.
                                      \label{eq:horizontalidentities76}
\end{align}

These tangent equations can be realized by genuine dilations.
Let $G_1(s)=G+sH_0$, $G_0(s)=I_d-G-sH_0$ for $s$ in a
sufficiently small open interval, and put $R_y(s)=\sqrt{G_y(s)}$.
The first line of \eqref{eq:horizontalidentities76} implies that
\[
 T_y=(B_y-R_y'(0))R_y(0)^{-1}
\]
is anti-Hermitian.  Indeed, multiplying $T_y+T_y^*$ on the left
and right by $R_y(0)$ gives zero by differentiating
$R_y(s)^2=G_y(s)$.  Thus
$\widetilde A_y(s)=\exp(sT_y)R_y(s)$ has derivative $B_y$ at
zero and satisfies
$\widetilde A_y(s)^*\widetilde A_y(s)=G_y(s)$ exactly.  A measurement
isometry is
\[
 W(s)\psi=\sum_{y=0}^1
              |y\rangle_C|y\rangle_L\widetilde A_y(s)\psi.
\]
Here $C$ is the observed outcome, and $L$ and the factor space of
$\widetilde A_y$ are environments.  Tracing the environments
implements $\mathcal M_{G(s)}$, also on inputs entangled with a
reference.  At zero, \eqref{eq:horizontalidentities76} gives
\begin{equation}\label{eq:horizontaloverlap76}
 W^*\dot W=0,\qquad \dot W^*\dot W=K^2.
\end{equation}

Fix a common tester with $N$ slots.  Purify its initial state and
all intervening operations, and retain the measurement environments
only for this comparison.  The derivative of its final purification
in the direction $H_0$ is the sum of $N$ terms, each having one
differentiated occurrence of $W$.  The terms are mutually orthogonal.
For two different differentiated slots, cancel the common isometries
after the later slot; the remaining overlap contains
$W^*\dot W\otimes I=0$.  Each term has norm at most
$\|\dot W\|_{\rm op}=\|K\|_{\rm op}$, independently of the
reference and memory dimensions.  Hence the derivative of the
purification has norm at most $\sqrt N\|K\|_{\rm op}$.
The pure-state trace-norm derivative bound, followed by contraction
under the partial trace, bounds the observed output derivative by
\begin{equation}\label{eq:horizontaltangent76}
 2\sqrt N\|K\|_{\rm op}.
\end{equation}
A publicly stopped tester is padded to $N$ slots by calls on fixed
dummy inputs, whose outputs are discarded.  The same purification
argument then applies to its original output.  Classical feedback is
included by purifying the common controlled operations.

Now let $H$ be an arbitrary Hermitian tangent at $G$, and set
$\kappa=N^{-1/2}$.  There is a unique Hermitian solution of
\begin{equation}\label{eq:regularizedsylvester76}
 SK+KS+\kappa K=H.
\end{equation}
Write $H=H_0+R$, with $H_0=SK+KS$ and $R=\kappa K$.
Both are legitimate tangents at the interior effect $G$.  The
derivative of a fixed tester's output is linear in the inserted
channel tangent.  The part arising from $H_0$ is bounded by
\eqref{eq:horizontaltangent76}.  The part arising from $R$ is
bounded by $2N\|R\|_{\rm op}$.  To verify the latter bound,
the one-call tangent has two opposite reference blocks
\[
 B_\rho=\tr_A[(R\otimes I)\rho],\qquad -B_\rho.
\]
Duality against Hermitian contractions gives
$\|B_\rho\|_1\le\|R\|_{\rm op}$ for every input state.
Insert this tangent into each of the $N$ slots in turn; all
subsequent common channels contract the Hermitian trace norm.
This proves the bound for its sum.  Thus the actual tangent $H$
has output derivative bounded by
\begin{equation}\label{eq:regularizedtangent76}
 2\sqrt N\|K\|_{\rm op}+2N\|R\|_{\rm op}
                         =4\sqrt N\|K\|_{\rm op}.
\end{equation}
The splitting is used only for the differential at an interior
effect; it does not postulate a physical mixture realizing the
entire path.

In an eigenbasis of $G$, with $v_i=g_i(1-g_i)$, equation
\eqref{eq:regularizedsylvester76} reads
\[
 K_{ij}=\frac{H_{ij}}{\sqrt{v_i}+\sqrt{v_j}+\kappa}.
\]
Since $(\sqrt{v_i}+\sqrt{v_j}+\kappa)^2\ge v_i+v_j+\tau$,
\begin{equation}\label{eq:regularizedformbound76}
 \sqrt N\|K\|_{\rm op}
 \le\sqrt N\|K\|_{\rm HS}
 \le g_{N,G}(H).
\end{equation}
For an interior path, integrate the resulting output derivative
bound for each fixed tester, and then take the supremum.  This
proves \eqref{eq:matrixpath76} without exchanging a supremum
with a derivative.

For a path in the closed body, apply the result to
$G_\epsilon(t)=(1-2\epsilon)G(t)+\epsilon I_d$, where
$0<\epsilon<1/2$.  At fixed $N$, the form
$g_{N,G}(H)$ is continuous on the whole body and bounded above
by $N\|H\|_{\rm HS}$.  The integrals therefore converge as
$\epsilon\downarrow0$.  The hybrid inequality
$d_N(\mathcal M_E,\mathcal M_F)\le2N\|E-F\|_{\rm op}$
gives continuity of the distances at both endpoints.  Passing to
the limit proves the asserted closed-body statement.
\end{proof}

\subsection{The midpoint comparison}

\begin{proof}[Proof of Theorem~\ref{thm:matrixmetric76}]
For the upper bound use $G(t)=(1-t)E+tF$.  The map
$G\mapsto V_G=G-G^2$ is operator concave.  For $0<t\le1/2$,
write $G(t)=(1-2t)E+2tM$ and discard the positive multiple of
$V_E$ in the concavity inequality.  Applying the same argument
from the other endpoint gives
\begin{equation}\label{eq:midpointvariance76}
 V_{G(t)}\succeq a(t)V_M,\qquad
             a(t)=2\min\{t,1-t\}.
\end{equation}
Left and right multiplication preserve this order as quadratic
forms on Hermitian matrices.  Since $a(t)\le1$,
$\mathcal A_{G(t)}\succeq a(t)\mathcal A_M$.
Inverting positive operators reverses their order, whence
\[
 g_{N,G(t)}(H)\le a(t)^{-1/2}Q_N(E,F).
\]
The scalar factor has integral $2$ on $[0,1]$.
Lemma~\ref{lem:matrixpath76} yields $d_N\le8Q_N$;
the trace-norm bound $d_N\le2$ gives the stated upper estimate.

For the lower bound choose any orthonormal eigenbasis $e_i$ of
$M$ and put
\begin{equation}\label{eq:matrixentries76}
 q_{ij}=|H_{ij}|\sqrt{\frac N{v_i+v_j+\tau}},
                 \qquad Q_N^2=\sum_{i,j}q_{ij}^2.
\end{equation}
Repeated input $e_i$ gives Bernoulli probabilities
$s=m_i-H_{ii}/2$, $t=m_i+H_{ii}/2$.  Their variances satisfy
\[
 \max\{s(1-s),t(1-t)\}
 \le s(1-s)+t(1-t)\le2m_i(1-m_i).
\]
Lemma~\ref{lem:bernoulliproduct75} therefore gives
\begin{equation}\label{eq:matrixdiagonallower76}
 d_N^{\rm na}(\mathcal M_E,\mathcal M_F)
                            \ge\tfrac1{64}\min\{1,q_{ii}\}.
\end{equation}

For $i\ne j$, restrict each input to
$\mathcal S=\spanop\{e_i,e_j\}$.  The resulting channels have
the compressed effects $E_{\mathcal S},F_{\mathcal S}$, including
on inputs entangled with a retained reference.  Their midpoint is
$M_{\mathcal S}=\diag(m_i,m_j)$.  Write their spectral data as
$(p,q,u)$ and $(p',q',v)$ in the notation of
Section~\ref{sec:biasedmetric75}, with gaps $r,r'$ and angle $h$
when both gaps are positive.  Put
\[
 V_E^{\mathcal S}=p(1-p)+q(1-q),\qquad
 V_F^{\mathcal S}=p'(1-p')+q'(1-q').
\]
For the compressed matrices the exact identity
\[
 V_{M_{\mathcal S}}=\tfrac12V_{E_{\mathcal S}}
              +\tfrac12V_{F_{\mathcal S}}
              +\tfrac14(F_{\mathcal S}-E_{\mathcal S})^2
\]
shows, on taking traces, that
\begin{equation}\label{eq:compressionvariance76}
 v_i+v_j+\tau\ge\tfrac12(V_E^{\mathcal S}+\tau),
 \qquad
 v_i+v_j+\tau\ge\tfrac12(V_F^{\mathcal S}+\tau).
\end{equation}
The scalar part of $F_{\mathcal S}-E_{\mathcal S}$ has zero
off-diagonal entry.  Its traceless part has operator norm
$|ru-r'v|/2$.  Decomposing this vector along the smaller radius
gives
\begin{align}
 |H_{ij}|&\le\tfrac12|r-r'|+\min\{r,r'\}\sin(h/2)\notag\\
 &\le\tfrac12\bigl(|p-p'|+|q-q'|+\min\{r,r'\}h\bigr).
                                      \label{eq:compressionoffdiagonal76}
\end{align}
When one radius is zero, the angular term is zero and the same
inequality holds without choosing its direction.  Each endpoint
variance is bounded by the appropriate $V_E^{\mathcal S}$ or
$V_F^{\mathcal S}$.  Equations
\eqref{eq:compressionvariance76}--\eqref{eq:compressionoffdiagonal76}
therefore imply
\begin{equation}\label{eq:compressionmodulus76}
 q_{ij}\le\frac1{\sqrt2}
 \left(T_N(p,p')+T_N(q,q')
       +\min\{K_N(p,q),K_N(p',q')\}h\right).
\end{equation}
The nonadaptive assertion of Theorem~\ref{thm:biasedmetric75}
applied to this compressed pair yields
\begin{equation}\label{eq:matrixoffdiagonallower76}
 d_N^{\rm na}(\mathcal M_E,\mathcal M_F)
                         \ge\tfrac1{8192}\min\{1,q_{ij}\}.
\end{equation}
Every such witness is a legitimate experiment for the original
effects by a fixed isometric embedding of $\mathcal S$.
Since $\max_{i,j}q_{ij}\ge Q_N/d$, taking the strongest of
\eqref{eq:matrixdiagonallower76} and
\eqref{eq:matrixoffdiagonallower76} proves the lower bound.
For $d=1$ only the diagonal argument is needed.

Finally, a one-use reference-assisted difference has output blocks
$B_\rho$ and $-B_\rho$, with
$B_\rho=\tr_A[((E-F)\otimes I)\rho]$.  Trace-norm duality
against Hermitian contractions gives
$\|B_\rho\|_1\le\|E-F\|_{\rm op}$.  An eigenstate of
$E-F$ at an extremal eigenvalue attains equality.  This proves
\eqref{eq:matrixoneuse76}.
\end{proof}

\begin{corollary}[Uniform comparison with nonadaptive experiments]
\label{cor:matrixadaptivity76}
For all $E,F\in\mathfrak E_d$ and all $N\ge1$,
\begin{equation}\label{eq:matrixadaptivity76}
 d_N(\mathcal M_E,\mathcal M_F)
           \le65536d\,d_N^{\rm na}(\mathcal M_E,\mathcal M_F).
\end{equation}
\end{corollary}
\begin{proof}
If $Q_N\le1$, divide the upper estimate $8Q_N$ by the lower
estimate $Q_N/(8192d)$.  If $Q_N\ge1$, the upper estimate is
$2$ and the lower estimate is $1/(8192d)$, which gives a smaller
constant.  The identical pair needs no division.
\end{proof}

\subsection{All spectral ranks and local forms}

\begin{lemma}[Spectral noncancellation in every rank]
\label{lem:matrixspectral76}
Let $\lambda_1(E)\ge\cdots\ge\lambda_d(E)$ and
$\lambda_1(F)\ge\cdots\ge\lambda_d(F)$ be the ordered spectra.
For every $1\le k\le d$,
\begin{equation}\label{eq:matrixspectrallower76}
 B_N(\lambda_k(E),\lambda_k(F))
                   \le d_N^{\rm na}(\mathcal M_E,\mathcal M_F).
\end{equation}
For effects diagonal in the same ordered orthonormal basis,
\begin{equation}\label{eq:matrixspectralupper76}
 d_N(\mathcal M_E,\mathcal M_F)
       \le\sum_{k=1}^d B_N(\lambda_k(E),\lambda_k(F)).
\end{equation}
\end{lemma}
\begin{proof}
Suppose $\lambda_k(E)\ge\lambda_k(F)$.  The span of the first
$k$ eigenvectors of $E$ intersects the orthogonal complement of
the first $k-1$ eigenvectors of $F$ in a nonzero subspace.  A unit
vector in that intersection has success probabilities
$s\ge\lambda_k(E)\ge\lambda_k(F)\ge t$ under the two
hypotheses.  The monotonicity in
Lemma~\ref{lem:bernoulliproduct75} implies
$B_N(s,t)\ge B_N(\lambda_k(E),\lambda_k(F))$.
Repeated use of this vector proves the first assertion; reverse
the hypotheses for the opposite ordering.  Multiplicities cause
no change in the dimension argument.

For the upper bound, first measure the common eigenbasis and then
produce a bit whose success probability is its corresponding
eigenvalue.  Supply $N$ independent bits in each of the $d$ rows
before the experiment.  Every common adaptive tester is a common
channel of these row arrays, including its reference blocks and
all publicly stopped histories.  The tensor-product triangle
inequality gives \eqref{eq:matrixspectralupper76}; unused row
bits are discarded.
\end{proof}

The following local comparison will permit volume estimates at
every boundary stratum.  It concerns matrices themselves and
does not require a coordinate chart for eigenvectors.

\begin{lemma}[Uniform local comparison of the variance forms]
\label{lem:matrixlocal76}
Set $W_G=G(I_d-G)+\tau I_d/2$.  For a Hermitian $H$, put
$q=g_{N,G}(H)$.  Then
\begin{equation}\label{eq:matrixrelative76}
 \|W_G^{-1/2}HW_G^{-1/2}\|_{\rm HS}\le2q.
\end{equation}
For either sign for which $G\pm H/2\in\mathfrak E_d$, one has
\begin{equation}\label{eq:matrixlocalperturbation76}
 \bigl\|W_G^{-1/2}(W_{G\pm H/2}-W_G)W_G^{-1/2}\bigr\|_{\rm op}
                       \le q+\tfrac34q^2.
\end{equation}
In particular, for $q\le1/4$,
\begin{equation}\label{eq:matrixlocalorder76}
 \tfrac12W_G\preceq W_{G\pm H/2}\preceq2W_G.
\end{equation}
Thus, if either $Q_N(E,F)\le1/4$ or
$g_{N,E}(F-E)\le1/4$, the midpoint and the fixed-endpoint
quadratic forms are comparable within factors $2$.
\end{lemma}
\begin{proof}
In an eigenbasis of $W_G$, with eigenvalues $w_i\ge\tau/2$,
\[
 \frac{w_i+w_j}{Nw_iw_j}\le4.
\]
Multiplying this inequality by
$N|H_{ij}|^2/(w_i+w_j)$ and summing proves
\eqref{eq:matrixrelative76}.  Put
$Z=W_G^{-1/2}HW_G^{-1/2}$.  The exact expansion
\[
 W_{G\pm H/2}-W_G
            =\pm\tfrac12(H-GH-HG)-\tfrac14H^2
\]
has normalized linear part
$\pm((I_d-2G)Z+Z(I_d-2G))/4$, since $G$ commutes with
$W_G$.  Its operator norm is at most $\|Z\|_{\rm op}/2\le q$.
The normalized quadratic part is $-ZW_GZ/4$.
Since $\tau\le1$, one has $\|W_G\|_{\rm op}\le3/4$, and
its norm is at most $3q^2/4$.  This proves
\eqref{eq:matrixlocalperturbation76} and hence
\eqref{eq:matrixlocalorder76}.  The same order holds for left
plus right multiplication, and inversion gives the quadratic-form
comparison.  Apply it first at $G=M$ or at $G=E$, respectively,
to obtain the final assertion.
\end{proof}
